\section{What constitutes evidence in an evolving ecosystem}
\label{sec:C01:evidence}
The literature and implementation review for all chapters extends through 29 October 2026; each source retains its reported revision and retrieval limits. Foundational methods are not excluded by age. Contemporary papers, official technical reports, normative specifications, and versioned repositories are examined according to the mechanism being claimed. First public release, manuscript revision, code release, and product activity are distinct dates. A 2026 revision of a 2025 paper is not automatically a newly introduced 2026 method. A university evaluation of a proprietary model is not research authored by that model's vendor.

\subsection{Search scope and inclusion rules}\label{sec:C01:search}
This is a mechanism-oriented critical survey, not a claim of exhaustive systematic identification. Candidate discovery combines a seed bibliography, citation following, and topic-specific searches for context editing, agent memory, procedural skills, tool protocols, harness optimization, agent evaluation, scientific agents, security, and serving. The present chapter additionally searches model--harness adaptation, pretraining capability elicitation, and harnessed reinforcement learning. Paper titles and identifiers are checked against original records; implementation claims use official repositories or documentation, with a commit, release, or protocol version when available. Retrieval occurred on 28--29 October 2026. The publication cutoff concerns artifact availability, not the calendar day on which an eligible artifact was inspected.

Inclusion requires a technical connection to an identified mechanism or evaluation boundary. Foundational references establish definitions, algorithms, or historical lineage; they do not substitute for current capability evidence. Frontier claims require an eligible primary artifact and retain its revision and experimental scope. A withdrawn manuscript does not support capability claims. Abstract-only or metadata-only inspection does not establish method details, and inaccessible implementation information remains an evidence gap. Organization-specific comparisons require verified paper affiliations rather than inferred ownership. Product catalogues provide ecosystem coverage, while feature and performance comparisons require additional mechanism-specific evidence.

For a mutable documentation page, the observation date and the available version have different meanings. A page read after the cutoff cannot establish that all of its present contents existed before the cutoff; a dated release or pinned artifact is needed for that historical claim. Consequently, the survey reports current mechanisms within the searched and inspected population rather than declaring any method universally best or all eligible work exhausted. Unresolved source leads are not counted as independently verified methods.

Four evidence classes are kept separate. An \emph{author-reported result} is attributed to an exact source version and experimental setting. An \emph{independent reproduction} requires actually executed artifacts, data identity, environment, configurations, seeds, and recorded outcomes. An \emph{analytical deduction} follows from stated assumptions and may provide a bound, counterexample, or resource calculation. A \emph{proposed experiment} defines a hypothesis and a discriminating intervention but has no observed result. This book does not present independent reproductions where none were performed.

Verification of a title and identifier does not verify a result. Inspection of a repository does not show that its current head reproduces a paper. A benchmark improvement does not prove the proposed mechanism caused the improvement unless resources and relevant alternatives are controlled. In particular, additional tool calls, longer reasoning, more candidates, or a stronger evaluator can account for gains attributed to a memory representation or orchestration strategy. The evidential question is the effect of the mechanism under declared operating conditions.

Conflicting findings are compared through task populations, resource accounting, feedback access, model versions, memory initialization, task order, and grading. Where these differences do not resolve the disagreement, the conflict remains open. A method can simultaneously improve average accuracy and worsen tail risk, or improve a selection score while losing held-out transfer. Those outcomes require a multi-dimensional interpretation rather than an editorial ranking. Evaluation frameworks such as HELM motivate systematic scenario and metric coverage, while agent benchmarks expose additional dependence on tool and environment semantics~\cite{liang2022helm,liu2023agentbench,yao2024taubench}.

\subsection{Contemporary measurements with distinct estimands}\label{sec:C01:currentevidence}
Agents' Last Exam evaluates configured agents on an executed workflow population; an expanding advertised task collection is not that experiment's denominator. Its mixed model--harness comparisons and fixed-harness backbone comparisons must remain distinct~\cite{corpus260605405}. AgentHorizon instead assesses judges on recorded instruction--trajectory pairs, including constructed near misses. It does not directly measure autonomous agent completion, and the constructed error distribution is not a natural deployment failure distribution~\cite{corpus261011050}. Table~\ref{tab:C01:currentevidence} makes the different inference boundaries explicit. All empirical findings remain author-reported.

{\small\setlength{\tabcolsep}{3pt}
\begin{table}[!htbp]
\centering
\caption{Current evidence anchors and limits of inference, inspected 29 October 2026.}\label{tab:C01:currentevidence}
\begin{tabular}{L{25mm}L{38mm}L{49mm}}
\toprule Artifact & Object measured or studied & Inference boundary\\\midrule
ALE, v4 & Configured-agent workflow outcomes. & Executed population and backbone/harness grouping define the comparison.\\
AgentHorizon, v1 & Judge detection of instruction violations in recorded trajectories. & Judge validity under the constructed distribution is distinct from agent reliability.\\
TraceTrial-Bench, v1 & Task outcomes and operational harness criteria in a selected execution archive. & Candidate failure links require evidence review; criterion satisfaction is not a causal defect estimate.\\
ASAP, v1 & Base-model execution under a completion scaffold, prefixes, and recovery. & Protocol-assisted elicitation does not identify unaided capability or later training gains.\\
Harness Scaling, v2 & Policy training and transfer across execution configurations. & Held-out combinations and external harness tests remain bounded by training and evaluation access.\\
\bottomrule
\end{tabular}
\end{table}
}

TraceTrial-Bench's selected archive is not a set of independent repetitions within every model--harness--task cell. Its semantic scoring also resolves some unexercised harness criteria using task outcomes, so the two evaluation layers are not statistically independent. Automatically activated mechanism links are candidates, not validated causal propagation~\cite{ch01_tracetrial2026}. ASAP's reference-prefix interventions and Harness Scaling's training exposure similarly belong in the experimental configuration~\cite{ch01_asap2026,ch01_harnessscaling2026}. These limitations change what a score estimates; they are not repaired by adding another decimal place or averaging unlike benchmarks.

\subsection{Provenance and population boundaries}\label{sec:C01:provenance}
Identity is a relation among dated artifacts, rather than a string copied from a product list. A renamed project, acquired service, and rewritten SDK can share a commercial lineage while exposing different execution semantics. An alias record should therefore specify the entities, the relation, its effective interval, and the primary evidence. A merger does not prove API compatibility, and an organization named in an announcement does not establish every author's affiliation. For an affiliation filter, the relevant observation is the author-to-institution statement in the inspected paper version; a repository owner or email domain is insufficient evidence of that paper's institutional authorship. These distinctions prevent organizational metadata from becoming an unsupported proxy for technical quality.

Release dates also have different meanings. The first preprint, a revised method, a public repository commit, a package release, and hosted availability are separate events. A claim about the earliest available mechanism must identify which event is relevant and retain the version that contained it. A later revision cannot be assigned the earlier version's date without checking that the mechanism was already present. Conversely, a preprint available before the literature cutoff remains eligible even when its stated conference occurs later. Dates establish availability, not independent validation or priority beyond the searched population.

Popularity measures need an equally explicit denominator. Repository stars, package downloads, public mentions, and service users measure different populations and have different duplication, access, and sampling effects. Record the platform, observation time, counting rule, entity identity, and missingness before comparing them. Such counts can describe visibility or adoption under that definition; they do not identify a capability ranking. A priority list can guide inspection without becoming evidence that its first entry is technically strongest.

An independent observatory supplies additional evidence only through a specified measurement protocol. Its task sample, evaluator, attempt budget, model endpoint, tool permissions, and observation date determine the estimand, just as they do for an author-run benchmark. Independence of organization does not guarantee independence of task data or judge errors. Agreement between two sites that reuse the same predictions is one result observed twice. Disagreement can instead arise from different configured systems or changing hosted versions. The appropriate synthesis preserves those boundaries and reports an unresolved comparison when they cannot be aligned.

